\documentclass[12pt,a4paper]{article}
\usepackage[utf8]{inputenc}
\usepackage[T5]{fontenc}
\usepackage[vietnamese]{babel}
\usepackage[margin=1.3cm,top=1.0cm,bottom=1.05cm,includehead,headheight=18pt,headsep=8pt]{geometry}
\usepackage{amsmath,amssymb,array,tabularx,enumitem,fancyhdr,tikz,xcolor,needspace}
\usetikzlibrary{calc,angles,quotes,arrows.meta,babel,patterns}
\definecolor{ink}{RGB}{30,55,70}

\newcommand{\headerQuote}{%
  \ifcase\value{page}\relax
    \or Math gives us hope that every problem has a solution% Trang 1
    \or Math + You = Champion% Trang 2
    \or Where Frustration Becomes Inspiration% Trang 3
    \or Turning Confusion into Clarity% Trang 4
    \or Don't be an X, be a Y of mathematics% Trang 5
    \or Don't get mad. Learn to add% Trang 6
    \or The only way to learn mathematics is to do mathematics% Trang 7
    \or Small steps every day lead to big results% Trang 8
    \or Every expert was once a beginner% Trang 9
    \or Success is the sum of small efforts repeated daily% Trang 10
    \or Mathematics is the music of reason% Trang 11
    \or Believe you can and you're halfway there% Trang 12
    \else Math gives us hope that every problem has a solution%
  \fi
}
\pagestyle{fancy}\fancyhf{}
\lhead{\small\textit{\headerQuote}}
\rhead{\small\textit{Người biên soạn: Dương Gia Bảo}}
\rfoot{\small Trang \thepage}
\renewcommand{\headrulewidth}{0.4pt}
\setlength{\headheight}{18pt}
\setlength{\parindent}{0pt}
\setlength{\parskip}{0pt}
\setlength{\tabcolsep}{1.5pt}
\renewcommand{\arraystretch}{1.0}
\newcommand{\parttitle}[1]{\par\vspace{1.5pt}\needspace{2\baselineskip}{\large\bfseries\color{ink}#1}\par\nopagebreak\vspace{1pt}}
\newcommand{\four}[4]{\par\noindent\begin{tabularx}{\linewidth}{@{}*{4}{>{\raggedright\arraybackslash}X}@{}}\textbf{A.} #1&\textbf{B.} #2&\textbf{C.} #3&\textbf{D.} #4\end{tabularx}}
\newcommand{\two}[4]{\par\noindent\begin{tabularx}{\linewidth}{@{}*{2}{>{\raggedright\arraybackslash}X}@{}}\textbf{A.} #1&\textbf{B.} #2\\[4.5pt]\textbf{C.} #3&\textbf{D.} #4\end{tabularx}}
\newcommand{\one}[4]{\par\noindent\begin{tabularx}{\linewidth}{@{}>{\raggedright\arraybackslash}X@{}}\textbf{A.} #1\\[4pt]\textbf{B.} #2\\[4pt]\textbf{C.} #3\\[4pt]\textbf{D.} #4\end{tabularx}}
\newcommand{\answerbox}{\hfill\textit{Đáp số:}\ \fbox{\rule{0pt}{9pt}\hspace{2.2cm}}}
\setlist[enumerate,1]{label=\textbf{Câu \arabic*.},leftmargin=1.1cm,labelsep=.1cm,itemsep=0.2pt,topsep=0.2pt,parsep=0pt}
\setlist[enumerate,2]{label=\alph*),leftmargin=.55cm,itemsep=1.5pt,topsep=0.2pt,parsep=0pt}

\begin{document}
\noindent
\begin{minipage}[c]{0.38\textwidth}
    \centering
    \textbf{LÀM TOÁN ĐỂ HẠNH PHÚC}\\[-2pt]
    \rule{3.2cm}{0.6pt}\\[3pt]
    \textbf{ĐỀ CHÍNH THỨC}
\end{minipage}\hfill
\begin{minipage}[c]{0.35\textwidth}
    \centering
    \textbf{ĐỀ ÔN THI GIỮA KÌ I}\\
    \textbf{NĂM HỌC 2026 -- 2027}\\[2pt]
    \textbf{ĐỀ SỐ 32}\\
    \textbf{Môn: TOÁN -- Khối lớp 10}\\
    \textit{\small Thời gian làm bài: 90 phút (không kể thời gian phát đề)}
\end{minipage}

\vspace{0.03cm}
\noindent
\begin{minipage}[b]{0.78\textwidth}
    \textbf{Họ và tên:}\dotfill\textbf{SBD:}\makebox[3.5cm]{\dotfill}
\end{minipage}\hfill
\begin{minipage}[b]{0.20\textwidth}
    \hfill
    \setlength{\fboxrule}{0.8pt}%
    \fbox{\begin{minipage}{2.3cm}
        \centering
        \textbf{\small Mã đề thi}\\
        \textbf{0032}
    \end{minipage}}
\end{minipage}
\vspace{0.03cm}

\parttitle{PHẦN I. Trắc nghiệm nhiều lựa chọn}
\textit{Thí sinh trả lời từ câu 1 đến câu 12. Mỗi câu hỏi thí sinh chỉ chọn một phương án.}
\begin{enumerate}
\item Cho hai mệnh đề $P$ và $Q$. Mệnh đề kéo theo $P \Rightarrow Q$ chỉ sai khi nào?
\two{$P$ đúng và $Q$ sai.}{$P$ sai và $Q$ đúng.}{Cả $P$ và $Q$ cùng sai.}{Cả $P$ và $Q$ cùng đúng.}

\item Cho mệnh đề $P\colon \text{``}\forall x \in \mathbb{R},\ x^2 - 3x + 5 > 0\text{''}$. Mệnh đề phủ định $\overline{P}$ của mệnh đề $P$ là
\two{$\overline{P}\colon \text{``}\forall x \in \mathbb{R},\ x^2 - 3x + 5 \le 0\text{''}$.}{$\overline{P}\colon \text{``}\exists x \in \mathbb{R},\ x^2 - 3x + 5 \le 0\text{''}$.}{$\overline{P}\colon \text{``}\exists x \in \mathbb{R},\ x^2 - 3x + 5 < 0\text{''}$.}{$\overline{P}\colon \text{``}\forall x \in \mathbb{R},\ x^2 - 3x + 5 < 0\text{''}$.}

\item Cho hai tập hợp $A = \{x \in \mathbb{N} \mid (2x - 1)(x^2 - 4) = 0\}$ và $B = \{x \in \mathbb{Z} \mid -1 \le x \le 2\}$. Khẳng định nào sau đây đúng?
\four{$A \cap B = \varnothing$.}{$A \cup B = \{0; 1; 2\}$.}{$B \setminus A = \{-1; 0; 1\}$.}{$A \setminus B = \{2\}$.}

\item Cho hai tập hợp $A = [-3; 5)$ và $B = (1; 7]$. Tập hợp hiệu $A \setminus B$ bằng
\four{$[-3; 1)$.}{$(1; 5)$.}{$(5; 7]$.}{$[-3; 1]$.}

\item Đường thẳng bờ của miền nghiệm bất phương trình bậc nhất hai ẩn $3x - 2y + 6 \ge 0$ là đường thẳng có phương trình nào sau đây?
\four{$3x - 2y + 6 = 0$.}{$3x + 2y + 6 = 0$.}{$2x - 3y + 6 = 0$.}{$3x - 2y - 6 = 0$.}

\item Cặp số $(x; y)$ nào sau đây là một nghiệm của bất phương trình bậc nhất hai ẩn $2x - y + 4 > 0$?
\four{$(1; 7)$.}{$(2; 3)$.}{$(-3; 0)$.}{$(0; 5)$.}

\item Đường thẳng $d$ đi qua hai điểm $A(0; 3)$ và $B(2; 0)$ có phương trình $3x + 2y - 6 = 0$. Nửa mặt phẳng kể cả bờ $d$ và chứa gốc toạ độ $O(0; 0)$ là miền nghiệm của bất phương trình nào sau đây?
\four{$3x + 2y - 6 > 0$.}{$3x - 2y - 6 \le 0$.}{$3x + 2y - 6 \le 0$.}{$3x + 2y - 6 \ge 0$.}

\item Cho tam giác $OAB$ với $O(0; 0)$, $A(3; 0)$, $B(0; 4)$. Tìm tất cả các giá trị của tham số $m$ để bất phương trình $x + 2y \le m$ nghiệm đúng với mọi điểm $(x; y)$ thuộc miền tam giác $OAB$ (kể cả biên).
\four{$m \ge 3$.}{$m \le 8$.}{$m \ge 7$.}{$m \ge 8$.}

\item Với mọi góc $\alpha$ thoả mãn $0^\circ < \alpha < 180^\circ$, khẳng định nào sau đây là đúng?
\two{$\sin(180^\circ - \alpha) = \sin \alpha$.}{$\cos(180^\circ - \alpha) = \cos \alpha$.}{$\tan(180^\circ - \alpha) = \tan \alpha$ ($\alpha \ne 90^\circ$).}{$\cot(180^\circ - \alpha) = \cot \alpha$.}

\item Giá trị của $\cos 120^\circ$ bằng
\four{$\dfrac{1}{2}$.}{$-\dfrac{1}{2}$.}{$-\dfrac{\sqrt{3}}{2}$.}{$\dfrac{\sqrt{3}}{2}$.}

\item Cho góc $\alpha$ thoả mãn $\tan \alpha = 3$. Giá trị của biểu thức $P = \dfrac{2\sin \alpha - 3\cos \alpha}{4\sin \alpha + \cos \alpha}$ bằng
\four{$\dfrac{5}{13}$.}{$-\dfrac{3}{13}$.}{$\dfrac{3}{13}$.}{$\dfrac{1}{3}$.}

\item Tam giác $ABC$ có độ dài hai cạnh $b = 8, c = 5$ và góc $\widehat{A} = 60^\circ$. Bán kính đường tròn ngoại tiếp $R$ của tam giác $ABC$ bằng
\four{$\dfrac{7}{2}$.}{$7\sqrt{3}$.}{$\dfrac{\sqrt{3}}{7}$.}{$\dfrac{7\sqrt{3}}{3}$.}
\end{enumerate}
\newpage

\parttitle{PHẦN II. Trắc nghiệm đúng/sai}
\textit{Thí sinh trả lời từ câu 1 đến câu 2. Trong mỗi ý a), b), c), d) ở mỗi câu, thí sinh chọn đúng hoặc sai.}
\begin{enumerate}
\item Cho hai tập hợp $A = \{x \in \mathbb{R} \mid x^2 - 4x + 3 = 0\}$, $B = \{x \in \mathbb{Z} \mid -1 \le x \le 3\}$ và tập hợp $C = [m - 2; m + 1]$ (với $m$ là tham số thực).
\begin{enumerate}
    \item Số phần tử của tập hợp $A$ bằng $2$.
    \item Tập hợp $B$ có tất cả $16$ tập hợp con.
    \item Giao của hai tập hợp $A$ và $B$ là $A \cap B = \{1; 3\}$.
    \item Có tất cả $6$ giá trị nguyên của tham số $m$ để $A \cap C \ne \varnothing$.
\end{enumerate}

\item Một cơ sở sản xuất bánh ngọt làm hai loại bánh: Bánh nướng và Bánh dẻo.
\begin{itemize}[leftmargin=1.5em, itemsep=0.1pt]
    \item Để làm một chiếc bánh nướng cần $40\text{ g}$ bột mì và $20\text{ g}$ đường.
    \item Để làm một chiếc bánh dẻo cần $30\text{ g}$ bột mì và $30\text{ g}$ đường.
\end{itemize}
Cơ sở hiện có tối đa $2{,}4\text{ kg}$ bột mì và $1{,}8\text{ kg}$ đường. Gọi $x$ và $y$ ($x, y \in \mathbb{N}$) lần lượt là số lượng bánh nướng và bánh dẻo mà cơ sở sản xuất.
\begin{enumerate}
    \item Số gam bột mì cần dùng để sản xuất $x$ chiếc bánh nướng và $y$ chiếc bánh dẻo là $40x + 30y$ (gam).
    \item Hệ bất phương trình ràng buộc số lượng bánh $x$ và $y$ có thể sản xuất là $\begin{cases} x \ge 0,\ y \ge 0 \\ 4x + 3y \le 240 \\ 2x + 3y \le 180. \end{cases}$
    \item Nếu cơ sở sản xuất $30$ chiếc bánh nướng và $40$ chiếc bánh dẻo thì không vượt quá lượng nguyên liệu hiện có.
    \item Biết mỗi chiếc bánh nướng bán được tiền lãi $15$ nghìn đồng và mỗi chiếc bánh dẻo bán được tiền lãi $18$ nghìn đồng. Số tiền lãi lớn nhất mà cơ sở có thể thu được là $1200$ nghìn đồng.
\end{enumerate}
\end{enumerate}

\parttitle{PHẦN III. Trắc nghiệm trả lời ngắn}
\textit{Thí sinh trả lời từ câu 1 đến câu 4. Điền kết quả vào ô tương ứng.}
\begin{enumerate}
\item Cho hai khoảng $A = (-3; 2)$ và $B = (m - 1; m + 4)$ (với $m$ là tham số thực). Có bao nhiêu giá trị nguyên của tham số $m$ để $A \cap B$ là một khoảng có độ dài bằng $3$?
\answerbox

\item Cho góc $\alpha$ thoả mãn $90^\circ < \alpha < 180^\circ$ và $\sin \alpha = \dfrac{5}{13}$. Tính giá trị của biểu thức $P = 13\cos \alpha + 12\tan \alpha$.
\answerbox
\end{enumerate}
\newpage

\begin{enumerate}[resume]
\item Một công ty dự định sản xuất hai loại sản phẩm $A$ và $B$ từ ba loại nguyên liệu $I, II, III$. Để sản xuất một tấn sản phẩm $A$ cần dùng $2$ tấn nguyên liệu $I$, $1$ tấn nguyên liệu $II$ và $1$ tấn nguyên liệu $III$. Để sản xuất một tấn sản phẩm $B$ cần dùng $1$ tấn nguyên liệu $I$, $1$ tấn nguyên liệu $II$ và $3$ tấn nguyên liệu $III$. Kho dự trữ của công ty hiện có tối đa $10$ tấn nguyên liệu $I$, $7$ tấn nguyên liệu $II$ và $15$ tấn nguyên liệu $III$. Mỗi tấn sản phẩm $A$ mang lại lợi nhuận $4$ triệu đồng, mỗi tấn sản phẩm $B$ mang lại lợi nhuận $5$ triệu đồng. Hỏi công ty có thể thu được lợi nhuận lớn nhất bằng bao nhiêu triệu đồng?
\answerbox

\item Một lớp học có $45$ học sinh. Trong đợt khảo sát sở thích, có $25$ học sinh thích môn Toán, $20$ học sinh thích môn Vật lí, $18$ học sinh thích môn Hoá học. Có $12$ học sinh thích cả Toán và Vật lí, $9$ học sinh thích cả Vật lí và Hoá học, $10$ học sinh thích cả Toán và Hoá học. Biết rằng có $5$ học sinh thích cả ba môn học trên. Hỏi có bao nhiêu học sinh trong lớp không thích môn nào trong ba môn Toán, Vật lí, Hoá học?
\answerbox
\end{enumerate}

\parttitle{PHẦN IV. Tự luận}
\textit{Thí sinh trình bày chi tiết lời giải từ câu 1 đến câu 4.}
\begin{enumerate}
\item Từ một trạm quan sát $O$ trên bờ biển, người ta theo dõi một tàu cá di chuyển trên biển. Tại thời điểm thứ nhất, tàu cá ở vị trí $A$ cách trạm quan sát $OA = 800\text{ m}$. Sau một khoảng thời gian, tàu cá di chuyển theo đường thẳng đến vị trí $B$ cách trạm quan sát $OB = 1200\text{ m}$. Góc nhìn giữa hai vị trí đo được là $\widehat{AOB} = 60^\circ$.
\begin{enumerate}
    \item Tính khoảng cách $AB$ giữa hai vị trí của tàu cá (làm tròn kết quả đến hàng đơn vị theo đơn vị mét).
    \item Trên đường đi thẳng từ vị trí $A$ đến vị trí $B$, hãy xác định khoảng cách ngắn nhất từ trạm quan sát $O$ đến con tàu (làm tròn kết quả đến hàng phần mười theo đơn vị mét).
\end{enumerate}

\item Một xe tải có tải trọng chở hàng tối đa là $5$ tấn và thể tích thùng xe chứa tối đa là $18\text{ m}^3$. Bác tài xế nhận chở hai loại kiện hàng $A$ và $B$:
\begin{itemize}[leftmargin=1.5em, itemsep=0.1pt]
    \item Mỗi kiện hàng loại $A$ nặng $0{,}2$ tấn và có thể tích $1\text{ m}^3$.
    \item Mỗi kiện hàng loại $B$ nặng $0{,}5$ tấn và có thể tích $1{,}2\text{ m}^3$.
\end{itemize}
\begin{enumerate}
    \item Gọi $x$ và $y$ ($x \ge 0, y \ge 0$) lần lượt là số kiện hàng loại $A$ và loại $B$ xếp lên xe tải. Hãy thiết lập hệ bất phương trình bậc nhất hai ẩn mô tả các ràng buộc về tải trọng và thể tích của xe tải.
    \item Nếu bác tài xế muốn chở $10$ kiện hàng loại $A$ và $6$ kiện hàng loại $B$ thì xe tải có chở được chuyến hàng này không? Vì sao?
\end{enumerate}

\item Cho hai tập hợp $A = [-3; 4)$ và $B = [m - 2; 2m + 1)$ với $m$ là tham số thực.
\begin{enumerate}
    \item Khi $m = 1$, hãy xác định các tập hợp $B$, $A \cap B$, $A \cup B$ và $A \setminus B$.
    \item Tìm tất cả các giá trị của tham số $m$ để tập hợp $B$ là một tập con khác rỗng của tập hợp $A$ ($B \ne \varnothing$ và $B \subset A$).
\end{enumerate}

\item Một hộ nông dân dự định đầu tư tối đa $60$ triệu đồng để nuôi hai loại gia súc: Bò và Lợn trên bãi đất của gia đình.
\begin{itemize}[leftmargin=1.5em, itemsep=0.1pt]
    \item Nuôi một con bò cần số vốn là $10$ triệu đồng và diện tích chuồng trại là $6\text{ m}^2$.
    \item Nuôi một con lợn cần số vốn là $2$ triệu đồng và diện tích chuồng trại là $2\text{ m}^2$.
\end{itemize}
Tổng diện tích chuồng trại có thể sử dụng tối đa là $48\text{ m}^2$. Sau một chu kì xuất chuồng, mỗi con bò mang lại lợi nhuận $4$ triệu đồng, mỗi con lợn mang lại lợi nhuận $1$ triệu đồng.
\begin{enumerate}
    \item Gọi $x, y$ lần lượt là số con bò và số con lợn được nuôi ($x \ge 0, y \ge 0$). Hãy lập hệ bất phương trình bậc nhất hai ẩn mô tả các ràng buộc về vốn đầu tư và diện tích chuồng trại.
    \item Biểu diễn miền nghiệm của hệ trên mặt phẳng toạ độ $Oxy$ và xác định số con bò, số con lợn hộ nông dân nên nuôi để thu được số tiền lợi nhuận lớn nhất.
\end{enumerate}
\end{enumerate}

\newpage

% ==================== ĐÁP ÁN VÀ LỜI GIẢI CHI TIẾT ====================
\parttitle{BẢNG ĐÁP ÁN TỔNG HỢP}

\vspace{0.2cm}
\noindent\textbf{PHẦN I. Trắc nghiệm nhiều lựa chọn (12 câu)}
\vspace{0.15cm}

\begin{center}
\begin{tabular}{|c|c|c|c|c|c|c|c|c|c|c|c|}
\hline
\textbf{1} & \textbf{2} & \textbf{3} & \textbf{4} & \textbf{5} & \textbf{6} & \textbf{7} & \textbf{8} & \textbf{9} & \textbf{10} & \textbf{11} & \textbf{12} \\
\hline
\textbf{A} & \textbf{B} & \textbf{C} & \textbf{D} & \textbf{A} & \textbf{B} & \textbf{C} & \textbf{D} & \textbf{A} & \textbf{B} & \textbf{C} & \textbf{D} \\
\hline
\end{tabular}
\end{center}

\vspace{0.3cm}
\noindent\textbf{PHẦN II. Trắc nghiệm đúng/sai (2 câu)}
\vspace{0.15cm}

\begin{center}
\begin{tabular}{|c|c|c|c|c|}
\hline
\textbf{Câu} & \textbf{Lệnh hỏi a)} & \textbf{Lệnh hỏi b)} & \textbf{Lệnh hỏi c)} & \textbf{Lệnh hỏi d)} \\
\hline
\textbf{Câu 1} & \textbf{Đúng} & \textbf{Sai} & \textbf{Đúng} & \textbf{Đúng} \\
\hline
\textbf{Câu 2} & \textbf{Đúng} & \textbf{Đúng} & \textbf{Đúng} & \textbf{Sai} \\
\hline
\end{tabular}
\end{center}

\vspace{0.3cm}
\noindent\textbf{PHẦN III. Trắc nghiệm trả lời ngắn (4 câu)}
\vspace{0.15cm}

\begin{center}
\begin{tabular}{|c|c|c|c|c|}
\hline
\textbf{Câu} & \textbf{Câu 1} & \textbf{Câu 2} & \textbf{Câu 3} & \textbf{Câu 4} \\
\hline
\textbf{Đáp số} & $2$ & $-17$ & $32$ & $8$ \\
\hline
\end{tabular}
\end{center}

\vspace{0.4cm}
\parttitle{PHẦN IV. ĐÁP SỐ TỰ LUẬN}
\vspace{0.15cm}

\begin{enumerate}[leftmargin=1.5em, itemsep=6pt]
    \item[\textbf{Câu 1.}] 
    \begin{itemize}[leftmargin=1.2em, itemsep=2pt]
        \item[a)] $AB = 400\sqrt{7} \approx 1058\text{ m}$.
        \item[b)] $OH = \dfrac{1200\sqrt{21}}{7} \approx 785{,}6\text{ m}$.
    \end{itemize}

    \item[\textbf{Câu 2.}] 
    \begin{itemize}[leftmargin=1.2em, itemsep=2pt]
        \item[a)] $\begin{cases} x \ge 0,\ y \ge 0 \\ 2x + 5y \le 50 \\ 5x + 6y \le 90 \end{cases}$
        \item[b)] Xe tải chở được chuyến hàng này (vì cặp $(10; 6)$ thoả mãn cả hai ràng buộc: $2(10) + 5(6) = 50 \le 50$ và $5(10) + 6(6) = 86 \le 90$).
    \end{itemize}

    \item[\textbf{Câu 3.}] 
    \begin{itemize}[leftmargin=1.2em, itemsep=2pt]
        \item[a)] $B = [-1; 3)$, \quad $A \cap B = [-1; 3)$, \quad $A \cup B = [-3; 4)$, \quad $A \setminus B = [-3; -1) \cup [3; 4)$.
        \item[b)] $-1 \le m \le \dfrac{3}{2}$.
    \end{itemize}

    \item[\textbf{Câu 4.}] 
    \begin{itemize}[leftmargin=1.2em, itemsep=2pt]
        \item[a)] $\begin{cases} x \ge 0,\ y \ge 0 \\ 5x + y \le 30 \\ 3x + y \le 24 \end{cases}$
        \item[b)] Hộ nông dân nên nuôi $3$ con bò và $15$ con lợn để thu được lợi nhuận lớn nhất là $27$ triệu đồng.
    \end{itemize}
\end{enumerate}

\newpage

% ==================== HÌNH VẼ MIỀN NGHIỆM HỆ BẤT PHƯƠNG TRÌNH ====================
\parttitle{HÌNH VẼ MIỀN NGHIỆM HỆ BẤT PHƯƠNG TRÌNH}

\vspace{0.2cm}
\noindent\textbf{1. Miền nghiệm Câu 3 -- Phần III (Bài toán công ty sản xuất sản phẩm $A$ và $B$):}
\[
\text{Hệ bất phương trình ràng buộc: }
\begin{cases}
2x + y \le 10 \\
x + y \le 7 \\
x + 3y \le 15 \\
x \ge 0,\ y \ge 0
\end{cases}
\]

\begin{center}
\begin{tikzpicture}[x=1.15cm, y=1.15cm]
    % Lưới toạ độ phụ
    \draw[help lines, color=gray!25, step=1] (-0.8,-0.8) grid (7.5,7.5);
    
    % Trục tọa độ
    \draw[-{Latex[length=3mm]}, line width=1.1pt] (-0.8,0) -- (7.8,0) node[below=2pt, font=\normalsize\bfseries] {$x$};
    \draw[-{Latex[length=3mm]}, line width=1.1pt] (0,-0.8) -- (0,7.8) node[left=2pt, font=\normalsize\bfseries] {$y$};
    \node[below left=2pt, font=\normalsize\bfseries] at (0,0) {$O$};
    
    % Tô màu miền nghiệm tứ giác OABD
    \fill[ink!20] (0,0) -- (5,0) -- (3,4) -- (0,5) -- cycle;
    \draw[line width=1.4pt, ink] (0,0) -- (5,0) -- (3,4) -- (0,5) -- cycle;
    
    % Các đường thẳng biên
    % d1: 2x + y = 10
    \draw[line width=1.1pt, blue!80!black] (1.5,7) -- (5.5,-1) 
        node[pos=0.05, above right, font=\small\bfseries] {$d_1: 2x + y = 10$};
        
    % d2: x + y = 7
    \draw[line width=1.1pt, red!80!black] (-0.5,7.5) -- (7.5,-0.5) 
        node[pos=0.92, above right, font=\small\bfseries] {$d_2: x + y = 7$};
        
    % d3: x + 3y = 15
    \draw[line width=1.1pt, teal!80!black] (-0.6,5.2) -- (7.5,2.5) 
        node[pos=0.95, above right, font=\small\bfseries] {$d_3: x + 3y = 15$};

    % Các đường gióng toạ độ
    \draw[dashed, line width=0.8pt, gray!80] (3,0) node[below=2pt, font=\small\bfseries] {$3$} -- (3,4) -- (0,4) node[left=2pt, font=\small\bfseries] {$4$};
    \node[below=2pt, font=\small\bfseries] at (5,0) {$5$};
    \node[left=2pt, font=\small\bfseries] at (0,5) {$5$};

    % Các điểm đỉnh
    \fill[black] (0,0) circle (2.5pt);
    \fill[black] (5,0) circle (2.5pt) node[above right=1pt, font=\normalsize\bfseries] {$A(5;0)$};
    \fill[black] (3,4) circle (2.5pt) node[above right=1pt, font=\normalsize\bfseries] {$B(3;4)$};
    \fill[black] (0,5) circle (2.5pt) node[above right=1pt, font=\normalsize\bfseries] {$D(0;5)$};

    % Nhãn miền nghiệm
    \node[font=\bfseries\color{ink}] at (2.0,2.2) {Miền nghiệm $OABD$};
\end{tikzpicture}
\end{center}

\vspace{0.2cm}
\begin{itemize}[leftmargin=1.5em, itemsep=3pt]
    \item Miền nghiệm là tứ giác $OABD$ (kể cả biên) với toạ độ các đỉnh:
    \[
    O(0;0), \quad A(5;0), \quad B(3;4), \quad D(0;5).
    \]
    *(Ba đường thẳng $d_1, d_2, d_3$ đồng quy tại điểm $B(3;4)$).*
    \item Hàm mục tiêu lợi nhuận: $F(x, y) = 4x + 5y$ (triệu đồng).
    \item Giá trị tại các đỉnh: $F(O) = 0$; $F(A) = 20$; $F(D) = 25$; $F(B) = 32$.
    \item \textbf{Kết luận:} Lợi nhuận lớn nhất là $\mathbf{32}$ triệu đồng khi sản xuất $3$ tấn sản phẩm $A$ và $4$ tấn sản phẩm $B$.
\end{itemize}

\newpage

% ==================== MIỀN NGHIỆM CÂU 4 PHẦN IV ====================
\parttitle{HÌNH VẼ MIỀN NGHIỆM HỆ BẤT PHƯƠNG TRÌNH (Tiếp theo)}

\vspace{0.2cm}
\noindent\textbf{2. Miền nghiệm Câu 4 -- Phần IV (Bài toán hộ nông dân nuôi bò và lợn):}
\[
\text{Hệ bất phương trình ràng buộc: }
\begin{cases}
x \ge 0,\ y \ge 0 \\
5x + y \le 30 \\
3x + y \le 24
\end{cases}
\]

\begin{center}
\begin{tikzpicture}[x=0.95cm, y=0.36cm]
    % Lưới toạ độ phụ
    \draw[help lines, color=gray!25, xstep=1, ystep=4] (-0.8,-2) grid (8.5,32);
    
    % Trục tọa độ
    \draw[-{Latex[length=3mm]}, line width=1.1pt] (-0.8,0) -- (9.0,0) node[below=2pt, font=\normalsize\bfseries] {$x$};
    \draw[-{Latex[length=3mm]}, line width=1.1pt] (0,-2) -- (0,33) node[left=2pt, font=\normalsize\bfseries] {$y$};
    \node[below left=2pt, font=\normalsize\bfseries] at (0,0) {$O$};

    % Tô màu miền nghiệm tứ giác OABC
    \fill[ink!20] (0,0) -- (6,0) -- (3,15) -- (0,24) -- cycle;
    \draw[line width=1.4pt, ink] (0,0) -- (6,0) -- (3,15) -- (0,24) -- cycle;

    % d1: 5x + y = 30
    \draw[line width=1.1pt, blue!80!black] (0.5,27.5) -- (6.5,-2.5) 
        node[pos=0.08, above right, font=\small\bfseries] {$d_1: 5x + y = 30$};

    % d2: 3x + y = 24
    \draw[line width=1.1pt, red!80!black] (-0.5,25.5) -- (8.5,-1.5) 
        node[pos=0.92, above right, font=\small\bfseries] {$d_2: 3x + y = 24$};

    % Đường gióng toạ độ
    \draw[dashed, line width=0.8pt, gray!80] (3,0) node[below=2pt, font=\small\bfseries] {$3$} -- (3,15) -- (0,15) node[left=2pt, font=\small\bfseries] {$15$};
    \node[below=2pt, font=\small\bfseries] at (6,0) {$6$};
    \node[left=2pt, font=\small\bfseries] at (0,24) {$24$};

    % Các đỉnh
    \fill[black] (0,0) circle (2.5pt);
    \fill[black] (6,0) circle (2.5pt) node[above right=1pt, font=\normalsize\bfseries] {$A(6;0)$};
    \fill[black] (3,15) circle (2.5pt) node[above right=1pt, font=\normalsize\bfseries] {$B(3;15)$};
    \fill[black] (0,24) circle (2.5pt) node[above right=1pt, font=\normalsize\bfseries] {$C(0;24)$};

    % Nhãn miền nghiệm
    \node[font=\bfseries\color{ink}] at (2.2,8.0) {Miền nghiệm $OABC$};
\end{tikzpicture}
\end{center}

\vspace{0.2cm}
\begin{itemize}[leftmargin=1.5em, itemsep=3pt]
    \item Miền nghiệm là tứ giác $OABC$ (kể cả biên) với toạ độ các đỉnh:
    \[
    O(0;0), \quad A(6;0), \quad B(3;15), \quad C(0;24).
    \]
    \item Hàm lợi nhuận: $L(x, y) = 4x + y$ (triệu đồng).
    \item Giá trị tại các đỉnh: $L(O) = 0$; $L(A) = 24$; $L(C) = 24$; $L(B) = 27$.
    \item \textbf{Kết luận:} Lợi nhuận lớn nhất là $\mathbf{27}$ triệu đồng khi nuôi $\mathbf{3}$ con bò và $\mathbf{15}$ con lợn.
\end{itemize}

\end{document}
